\end{itemize}
\end{theorem}

Theorem~\ref{theorem2} coincides with the results in~\cite{LuMazoZhao}.

\subsection[Proof for $p \neq 0$]{Proof for $\boldsymbol{p \neq 0}$}
\label{PFqn0}

The Verma modules $ V^{\delta,p}_{\ell} $ are reducible for $ \ell \geq 2 $ so we restrict ourselves to $ \ell \geq 2$.
We consider the quotient module $ V^{\delta,p}_{\ell}/{\cal I}^{(2)} $ where $ {\cal I}^{(2)} $ is the largest
${\mathfrak g}_{\ell}$-submodule of $ V^{\delta,p}_{\ell}$.
Since there is no singular vectors in the $ N = 1 $ subspace $ (V^{\delta,p}_{\ell})_1$, $ {\cal I}^{(2)} $ will be
induced by the singular vector in the $ N = 2 $ subspace.
\begin{lemma}
\label{lemma3}
There exists precisely one $($up to an overall constant$)$ singular vector in the $ N = 2 $ subspace $
(V^{\delta,p}_{\ell})_2$.
This singular vector is given~by
\begin{gather*}
\big|v_s^{(2)}\big\rangle = (2p (\ell+1) P_{\ell-2} - (\ell+2) P_{\ell-1}^2) \ket{\delta,p}.
%\label{SVN2}
\end{gather*}
\end{lemma}


\begin{proof}
The basis of $ \big(V^{\delta,p}_{\ell}\big)_2 $ is $ \ket{2,\underaccent{\tilde}{0}}, \ket{1,\underaccent{\tilde}{\epsilon}_1},
\ket{0,2\underaccent{\tilde}{\epsilon}_1}, \ket{0,\underaccent{\tilde}{\epsilon}_2}$.
The singular vector $ \big|v_s^{(2)}\big\rangle$ is a~linear combination of the basis:
\begin{gather*}
\big|v_s^{(2)}\big\rangle  = c_1 \ket{2,\underaccent{\tilde}{0}} + c_2 \ket{1,\underaccent{\tilde}{\epsilon}_1} + c_3
\ket{0,2\underaccent{\tilde}{\epsilon}_1} + c_4 \ket{0,\underaccent{\tilde}{\epsilon}_2}.
\end{gather*}
It must satisfy the condition
\begin{gather*}
0 = P_{\ell+1} \big|v_s^{(2)}\big\rangle  = 2 (\ell+1)  p   c_1 \ket{1,\underaccent{\tilde}{0}} + (\ell+1)   (\ell  c_1 + p
c_2) \ket{0,\underaccent{\tilde}{\epsilon}_1}.
\end{gather*}
Thus $ c_1 = c_2 = 0$.
Furthermore, the condition $ C \big|v_s^{(2)}\big\rangle  = 0 $ yields the relation
\begin{gather*}
2 (\ell+1)   p   c_3 + (\ell+2) c_4 = 0.
\end{gather*}
This proves Lemma~\ref{lemma3}.
\end{proof}

Def\/ine $ {\cal I}^{(2)} = U({\mathfrak g}_{\ell}^{+}) \big|v_s^{(2)}\big\rangle $, then $ {\cal I}^{(2)} $ is the largest $
{\mathfrak g}_{\ell}$-submodule in $ V^{\delta,p}_{\ell}$.
Let $ \big|u_0^{(2)}\big\rangle  $ be the lowest weight vector in $ V^{\delta,p}_{\ell}/{\cal I}^{(2)}$.
Then
\begin{gather*}
    D \big|u_0^{(2)}\big\rangle  = \delta \big|u_0^{(2)}\big\rangle ,
\qquad
P_{\ell} \big|u_0^{(2)}\big\rangle  = p \big|u_0^{(2)}\big\rangle ,
\nonumber
\\
    X \big|u_0^{(2)}\big\rangle  = 0, \qquad X \in {\mathfrak g}_{\ell}^-,
\qquad
P_{\ell-2} \big|u_0^{(2)}\big\rangle  = \frac{\ell+2}{2p (\ell+1)} P_{\ell-1}^2 \big|u_0^{(2)}\big\rangle .
%\label{LWVQuo2}
\end{gather*}
It follows that the basis of $ V^{\delta,p}_{\ell}/{\cal I}^{(2)} $ is given~by
\begin{gather*}
\big|k,\underaccent{\tilde}{m}^{(2)}\big\rangle  = H^k   P_{\ell-1}^{m_1}  P_{\ell-3}^{m_3} \cdots P_{0}^{m_{\ell}}
\big|u_0^{(2)}\big\rangle ,
%\label{BasisQuo2}
\end{gather*}
where $ \underaccent{\tilde}{m}^{(2)} = (m_1, 0, m_3, \dots,m_{\ell}) \in {\mathbb R}^{\ell}$.
Observing the relation
\begin{gather*}
D \big|k,\underaccent{\tilde}{m}^{(2)}\big\rangle  = \left(\delta+k +m_1 + \sum\limits_{i=3}^{\ell}i m_i\right)
\big|k,\underaccent{\tilde}{m}^{(2)}\big\rangle ,
\end{gather*}
we def\/ine the level $ N^{(2)} $ in the quotient space $ V^{\delta,p}_{\ell}/{\cal I}^{(2)} $ by $  N^{(2)}
= k + m_1 + \sum\limits_{i=3}^{\ell} i m_{i}$.
The vectors $ \{
\ket{1,\underaccent{\tilde}{0}},
\ket{0,\underaccent{\tilde}{\epsilon}_1}
\} $ and $ \{
\ket{2,\underaccent{\tilde}{0}},
\ket{1,\underaccent{\tilde}{\epsilon}_1},
\ket{0,2\underaccent{\tilde}{\epsilon}_1}
\}$ form a~basis of $ N^{(2)} = 1 $ and $ N^{(2)} = 2 $ subspaces of $ V^{\delta,p}_{\ell}/{\cal I}^{(2)}$,
respectively.
The Shapovalov determinant for $ N^{(2)} = 1 $ subspace is same as~\eqref{SEC:VM} and given by $ (\ell+1)^2 p^2$.
While that for $ N^{(2)} = 2 $ subspace is calculated as follows
\begin{gather*}
    \left|
\begin{matrix} \bracket{2,\underaccent{\tilde}{0}}{2,\underaccent{\tilde}{0}} &
\bracket{2,\underaccent{\tilde}{0}}{1,\underaccent{\tilde}{\epsilon}_1} &
\bracket{2,\underaccent{\tilde}{0}}{0,2\underaccent{\tilde}{\epsilon}_1}
\vspace{1mm}\\ \bracket{1,\underaccent{\tilde}{\epsilon}_1}{2,\underaccent{\tilde}{0}} &
\bracket{1,\underaccent{\tilde}{\epsilon}_1}{1,\underaccent{\tilde}{\epsilon}_1} &
\bracket{1,\underaccent{\tilde}{\epsilon}_1}{}
\vspace{1mm}\\ \bracket{0,2\underaccent{\tilde}{\epsilon}_1}{2,\underaccent{\tilde}{0}} &
\bracket{0,2\underaccent{\tilde}{\epsilon}_1}{1,\underaccent{\tilde}{\epsilon}_1} &
\bracket{0,2\underaccent{\tilde}{\epsilon}_1}{0,2\underaccent{\tilde}{\epsilon}_1}
\end{matrix}
\right|
\\
\qquad{}
= \left|
\begin{matrix} 4(2\delta+1) \delta & 2(2\delta+1)(\ell+1)p & 2(\ell+1)^2 p^2
\vspace{1mm}\\ 2(2\delta+1)(\ell+1)p & (\ell+1)^2 p^2 & 0
\vspace{1mm}\\ 2(\ell+1)^2 p^2 & 0 & 0
\end{matrix}
\right| = -4(\ell+1)^6 p^6.
\end{gather*}
Therefore there exist no singular vectors in $ N^{(2)} = 1, 2 $ subspaces.
On the other hand, one f\/inds a~singular vector in the $ N^{(2)} = 3 $ subspace.

\begin{lemma}
\label{lemma4}
There exists precisely one $($up to overall constant$)$ singular vector in the $ N^{(2)} = 3 $ subspace of $
V^{\delta,p}_{\ell}/{\cal I}^{(2)}$.
The singular vector is given~by
\begin{gather*}
\big|v_s^{(3)}\big\rangle  = \big(3! p^2 (\ell+1)^2  P_{\ell-3} - (\ell+2)  (\ell+3)  P_{\ell-1}^3\big) \big|u_0^{(2)}\big\rangle .
%\label{SVN3}
\end{gather*}
\end{lemma}

\begin{proof}
The lemma can be proved in a~way exactly similar to Lemma~\ref{lemma3}.
The basis of the level $ N^{(2)} = 3 $ subspace is given~by
\begin{gather*}
\ket{3,\underaccent{\tilde}{0}},
\qquad
\ket{2,\underaccent{\tilde}{\epsilon}_1},
\qquad
\ket{1,2\underaccent{\tilde}{\epsilon}_1},
\qquad
\ket{0,3\underaccent{\tilde}{\epsilon}_1},
\qquad
\ket{0,\underaccent{\tilde}{\epsilon}_3}.
\end{gather*}
The singular vector $ \big|v_s^{(3)}\big\rangle  $ is a~linear combination of the basis:
\begin{gather*}
\big|v_s^{(3)}\big\rangle  = c_1 \ket{3,\underaccent{\tilde}{0}} + c_2 \ket{2,\underaccent{\tilde}{\epsilon}_1} + c_3
\ket{1,2\underaccent{\tilde}{\epsilon}_1} + c_4 \ket{0,3\underaccent{\tilde}{\epsilon}_1} + c_5
\ket{0,\underaccent{\tilde}{\epsilon}_3}.
\end{gather*}
It must satisfy the condition
\begin{gather*}
0  =  P_{\ell+1} \big|v_s^{(3)}\big\rangle
\\
\hphantom{0}{}
 =  3 p c_1 \ket{2,\underaccent{\tilde}{0}} + (3\ell c_1 + 2p c_2) \ket{1,\underaccent{\tilde}{\epsilon}_1} +
\left(\frac{(\ell+2) \ell (\ell-1)}{2(\ell+1)p} c_1 + \ell c_2 + (\ell+1) p c_3 \right) \ket{0,
2\underaccent{\tilde}{\epsilon}_1}.
\end{gather*}
Thus $ c_1 = c_2 = c_3 = 0$.
Furthermore, the condition $ C \big|v_s^{(3)}\big\rangle  = 0 $ yields the relation
\begin{gather*}
3 (\ell+1) p c_4 + \frac{(\ell+3) (\ell+2)}{2(\ell+1)p} c_5 = 0.
\end{gather*}
This proves  Lemma~\ref{lemma4}.
\end{proof}

The vector $ \big|v_s^{(3)}\big\rangle  $ is singular only in the quotient space $ V^{\delta,p}_{\ell}/{\cal I}^{(2)} $ and not in
$ V^{\delta,p}_{\ell} $ itself.
Such vector is called \textit{subsingular}~\cite{Dob95,Dob97}.

It follows from Lemma~\ref{lemma4} that the subspace $ {\cal I}^{(3)} = U({\mathfrak g}_{\ell}^+) \big|v_s^{(3)}\big\rangle  $ is
the largest ${\mathfrak g}_{\ell}$-submodule in $ V^{\delta,p}_{\ell}/{\cal I}^{(2)}$.
Now we consider the quotient space $ V^{\delta,p}_{\ell}/{\cal I}^{(2)}/{\cal I}^{(3)}:= \big(V^{\delta,p}_{\ell}/{\cal
I}^{(2)}\big)/{\cal I}^{(3)}$.
The lowest weight vector $ \big|u_0^{(3)}\big\rangle  $ of this quotient space is def\/ined~by
\begin{gather*}
    D \big|u_0^{(3)}\big\rangle  = \delta \big|u_0^{(3)}\big\rangle ,
\qquad
P_{\ell} \big|u_0^{(3)}\big\rangle  = p \big|u_0^{(3)}\big\rangle ,
\qquad
    X \big|u_0^{(3)}\big\rangle  = 0, \qquad X \in {\mathfrak g}_{\ell}^-,
\nonumber
\\
    P_{\ell-2} \big|u_0^{(3)}\big\rangle  = \frac{\ell+2}{2p  (\ell+1)} P_{\ell-1}^2 \big|u_0^{(3)}\big\rangle ,
\qquad
P_{\ell-3} \big|u_0^{(3)}\big\rangle  = \frac{(\ell+2) (\ell+3)}{3!  p^2  (\ell+1)^2} P_{\ell-1}^3\big|u_0^{(3)}\big\rangle .
%\label{LWVQuo3}
\end{gather*}
The basis of $ V^{\delta,p}_{\ell}/{\cal I}^{(2)}/{\cal I}^{(3)} $ is given~by
\begin{gather*}
H^k  P_{\ell-1}^{m_1} P_{\ell-4}^{m_4} \cdots P_0^{m_{\ell}} \big|u_0^{(3)}\big\rangle ,
\end{gather*}
and we def\/ine the level $  N^{(3)} = k + m_1 + \sum\limits_{i=4}^{\ell} i  m_i$.
With this setting one can show the followings:
\begin{enumerate}\itemsep=0pt
\item[i)] the Shapovalov determinant does not vanish for $ 1 \leq N^{(3)} \leq 3 $ so that there exist no singular vectors
in the subspaces with $ N^{(3)} = 1, 2, 3$;
\item[ii)] there exists a~unique singular vector in the level $ N^{(3)} = 4 $ subspace.
\end{enumerate}
By the singular vector in $ N^{(3)} = 4 $ subspace, one can def\/ine the largest ${\mathfrak g}_{\ell}$-submodule $ {\cal
I}^{(4)} $ in $ V^{\delta,p}_{\ell}/{\cal I}^{(2)}/{\cal I}^{(3)} $ and consider the quotient space $
V^{\delta,p}_{\ell}/{\cal I}^{(2)}/{\cal I}^{(3)}/{\cal I}^{(4)}$.

In fact, one can repeat this process until we arrive at $ V^{\delta,p}_{\ell}/{\cal I}^{(2)}/\cdots/{\cal I}^{(\ell)}$.
This is assured by the next lemma:
\begin{lemma}
\label{lemma5}
Suppose that we have arrived at the quotient space $ V_{\ell}^{(\lambda)}:= V^{\delta,p}_{\ell}/{\cal
I}^{(2)}/\cdots/{\cal I}^{(\lambda)}$, $2 \leq \lambda \leq \ell$.
Namely, we have the quotient space with the basis
\begin{gather*}
H^k P_{\ell-1}^{m_1} P_{\ell-\lambda-1}^{m_{\lambda+1}} \cdots P_0^{m_{\ell}} \big|u_0^{(\lambda)}\big\rangle ,
%\label{LWVQuoL}
\end{gather*}
where $ \big|u_0^{(\lambda)}\big\rangle  $ is the lowest weight vector in $ V_{\ell}^{(\lambda)} $, defined~by
\begin{gather}
    D \big|u_0^{(\lambda)}\big\rangle  = \delta \big|u_0^{(\lambda)}\big\rangle ,
\qquad
P_{\ell} \big|u_0^{(\lambda)}\big\rangle  = p \big|u_0^{(\lambda)}\big\rangle ,
\qquad
X \big|u_0^{(\lambda)}\big\rangle  = 0, \qquad X \in {\mathfrak g}_{\ell}^-,
\nonumber
\\      P_{\ell-a} \big|u_0^{(\lambda)}\big\rangle  = \frac{(\ell+2) (\ell+3) \cdots (\ell+a)}{a!   p^{a-1}   (\ell+1)^{a-1}}
P_{\ell-1}^a \big|u_0^{(\lambda)}\big\rangle ,
\qquad
a = 2, 3, \dots, \lambda.
\label{Assumption-lambda}
\end{gather}
Then
\begin{enumerate}\itemsep=0pt
\item[$i)$] $ V_{\ell}^{(\lambda)} $ is the graded vector space:
\begin{gather*}
    V_{\ell}^{(\lambda)} = \bigoplus_{N^{(\lambda)}=0}^{\infty} (V_{\ell}^{(\lambda)})_{N^{(\lambda)}},
\qquad
\big(V_{\ell}^{(\lambda)}\big)_{N^{(\lambda)}} = \big\{ \ket{v} \in V_{\ell}^{(\lambda)} \, | \, D \ket{v} = (\delta +
N^{(\lambda)}) \ket{v} \big\},
\nonumber
\\
    N^{(\lambda)} = k + m_1 + \sum\limits_{i=\lambda+1}^{\ell} i m_i.
%\label{gradedquotientspace}
\end{gather*}
\item[$ii)$] The subspace $ \big(V_{\ell}^{(\lambda)}\big)_{N^{(\lambda)}} $ has a~nonvanishing Shapovalov determinant if $ 1 \leq
N^{(\lambda)} \leq \lambda$.
The Shapovalov determinant is given by $($up to a~sign factor$)$
\begin{gather}
\Delta_{N^{(\lambda)}} = (p  (\ell+1))^{N^{(\lambda)}(N^{(\lambda)}+1)} \prod\limits_{k=0}^{N^{(\lambda)}} k!
\big(N^{(\lambda)}-k\big)!.
\label{KDinQuotient}
\end{gather}
This implies that there exists no singular vectors in the level $ N^{(\lambda)} $ subspaces if $ 1 {\leq} N^{(\lambda)}
{\leq} \lambda$.
\item[$iii)$] If $ 2 \leq \lambda \leq \ell-1$, then there exists precisely one $($up to overall constant$)$ singular vector in the
$ N^{(\lambda)} = \lambda + 1 $ subspace of $ V^{(\lambda)}_{\ell}$.
The singular vector is given~by
\begin{gather}
\big|v_s^{(\lambda{+}1)}\big\rangle  = \big(  (\lambda+1)!  p^{\lambda} (\ell+1)^{\lambda} P_{\ell{-}\lambda{-}1} \!- (\ell+2) (\ell+3)
\cdots (\ell+\lambda+1) P_{\ell{-}1}^{\lambda{+}1}  \big)\! \big|u_0^{(\lambda)}\big\rangle .\!\!\!\!
\label{SVinQuotientSpace}
\end{gather}
\end{enumerate}
\end{lemma}

\begin{proof}
(i) Can be trivially proved.
 (ii) Suppose that $ 1 \leq N^{(\lambda)} \leq \lambda$, then the basis of $
\big(V_{\ell}^{(\lambda)}\big)_{N^{(\lambda)}} $ is given by $ \ket{k} = H^{N^{(\lambda)}-k} P_{\ell-1}^k \big|u_0^{(\lambda)}\big\rangle
$ with $ k = 0, 1, \dots, N^{(\lambda)}$.
The equality in the next equation is up to a~sign factor:
\begin{gather}
\Delta_{N^{(\lambda)}} = \left|
\begin{matrix} \bracket{0}{N^{(\lambda)}} & \bracket{0}{N^{(\lambda)}-1} & \dots & \bracket{0}{0}
\\ \bracket{1}{N^{(\lambda)}} & \bracket{1}{N^{(\lambda)}-1} & \dots & \bracket{1}{0}
\\
\vdots & \vdots & & \vdots
\\
\bracket{N^{(\lambda)}}{N^{(\lambda)}} & \bracket{N^{(\lambda)}}{N^{(\lambda)}-1} & \vdots & \bracket{N^{(\lambda)}}{0}
\end{matrix}
\right|.
\label{KDQuoSpa}
\end{gather}
The lower triangular entry of~\eqref{KDQuoSpa} is given~by
\begin{gather}
\big\langle m \big|  N^{(\lambda)}-k\big\rangle  = \big\langle u_0^{(\lambda)}\big| P_{\ell+1}^m C^{N^{(\lambda)}-m} H^k P_{\ell-1}^{N^{(\lambda)}-k}
\big|u_0^{(\lambda)}\big\rangle ,
\qquad
k < m.
\label{Prod-m-Nk}
\end{gather}
Using the commutation relations
\begin{gather*}
\big[C, H^k\big] = 2k H^{k-1} D + k(k-1) H^{k-1},
\qquad
\big[C, P_{\ell-1}^k\big] = k (\ell+1) P_{\ell-1}^{k-1} P_{\ell},
\end{gather*}
we see that~\eqref{Prod-m-Nk} is a~linear combination of
\begin{gather}
\big\langle u_0^{(\lambda)}\big| P_{\ell+1}^m H^j P_{\ell-1}^{m-j} \big|u_0^{(\lambda)}\big\rangle ,
\qquad
\max\{k-N+m, 0\} \leq j \leq k.
\label{Prod-m-Nk-2}
\end{gather}
The lower bound for~$j$ corresponds to the two possibilities $ N-m \leq k $ or $ N-m > k$.
By the relation
\begin{gather*}
\big[H, P_{\ell+1}^m\big] = -m(\ell+1) P_{\ell+1}^{m-1} P_{\ell},
\end{gather*}
the equation~\eqref{Prod-m-Nk-2} is further reduced to
\begin{gather*}
\big\langle u_0^{(\lambda)}\big| P_{\ell+1}^{m-j} P_{\ell-1}^{m-j} \big|u_0^{(\lambda)}\big\rangle  = 0.
\end{gather*}
The last equality stems from $ m -j > 0 $ which is due to the range of~$j$ and $ k < m$.
Thus all the lower triangular entries of~\eqref{KDQuoSpa} vanish.

The diagonal entries of~\eqref{KDQuoSpa} correspond to~\eqref{Prod-m-Nk} with $ m = k$.
Thus they are expanded into a~linear combination of~\eqref{Prod-m-Nk-2}.
Because of $ m = k $ there exist precisely one term in the expansion which gives the nonvanishing contribution.
Writing the coef\/f\/icient explicitly, that term is given~by
\begin{gather*}
(N^{(\lambda)}-k) ((\ell+1) p)^{N^{(\lambda)}-k} \big\langle u_0^{(\lambda)}\big| P_{\ell+1}^k H^k \big|u_0^{(\lambda)}\big\rangle  = k!
(N^{(\lambda)}-k)!   (p  (\ell+1))^{N^{(\lambda)}}.
\end{gather*}
The formula~\eqref{KDinQuotient} immediately follows from these results.

 (iii) The basis vectors of level $ N^{(\lambda)} = \lambda + 1 $ subspace are $ P_{\ell-\lambda-1}
\big|u_0^{(\lambda)}\big\rangle  $ and $ \ket{k} =$ \linebreak  $H^{\lambda+1-k}   P_{\ell-1}^k \big|u_0^{(\lambda)}\big\rangle$ with $
k=0,1,\dots,\lambda+1$.
The singular vector is a~linear combination of these vectors:
\begin{gather*}
\big|v_s^{(\lambda+1)}\big\rangle  = \sum\limits_{k=0}^{\lambda+1} \alpha_k \ket{k} + \beta P_{\ell-\lambda-1}
\big|u_0^{(\lambda)}\big\rangle .
\end{gather*}
The condition $ P_{\ell+1} \big|v_s^{(\lambda+1)}\big\rangle  = 0 $ yields
\begin{gather*}
    (\lambda+1-n) \left(  p (\ell+1)  \alpha_n + \binom{\ell+1}{2} (\lambda+2-n) \alpha_{n-1} \right) H^{\lambda-n}
P_{\ell-1}^n \big|u_0^{(\lambda)}\big\rangle
\\
\qquad{}
+ \sum\limits_{j=2}^n \binom{\ell+1}{j+1} \frac{(\lambda+1-n+j)!}{(\lambda-n)!}   \alpha_{n-j} H^{\lambda-n} P_{\ell-j}
P_{\ell-1}^{n-j} \big|u_0^{(\lambda)}\big\rangle  = 0,
\qquad
0 \leq n \leq \lambda.
\end{gather*}
It follows that $ \alpha_j = 0 $ for $ 0 \leq j \leq \lambda$.
Thus the singular vector yields
\begin{gather*}
\big|v_s^{(\lambda+1)}\big\rangle  = \big(  \alpha_{\lambda+1} P_{\ell-1}^{\lambda+1} + \beta P_{\ell-\lambda-1}  \big)
\big|u_0^{(\lambda)}\big\rangle .
\end{gather*}
The condition $ C \big|v_s^{(\lambda+1)}\big\rangle  = 0 $ gives the relation
\begin{gather*}
(\lambda+1)!  p^{\lambda} (\ell+1)^{\lambda} \alpha_{\lambda+1} + (\ell+2) (\ell+3) \cdots (\ell+\lambda+1) \beta = 0.
\end{gather*}
This implies the uniqueness of the formula~\eqref{SVinQuotientSpace} of the singular vector.
We remark that if $ \lambda=\ell $ then the vector corresponds to $ P_{\ell-\lambda-1} $ does not exist.
Thus $ \alpha_{\lambda+1} = 0 $, so that no singular vectors at level $ \ell+1$.
\end{proof}

Now we consider the module $ V^{(\ell)}_{\ell} = V^{\delta,p}_{\ell}/{\cal I}^{(2)}/ {\cdots } /{\cal I}^{(\ell)}$.
The basis of this space is $ H^k P_{\ell{-}1}^m \big|u_0^{(\ell)}\big\rangle  $ where $ \big|u_0^{(\ell)}\big\rangle  $ is the lowest weight
vector def\/ined by~\eqref{Assumption-lambda} with $ \lambda = \ell$.
By Lemma~\ref{lemma5}, there exist no singular vectors in the subspaces of $ V^{(\ell)}_{\ell} $ labelled by $
N^{(\ell)} = 1, 2, \dots, \ell$.
For $ N^{(\ell)} \geq \ell+1 $ the singular vectors may be written as
\begin{gather*}
\big|v_s^{(\ell)}\big\rangle  = \sum\limits_{k=0}^{N^{(\ell)}} \alpha_k H^k P_{\ell-1}^{N^{(\ell)}-k} \big|u_0^{(\ell)}\big\rangle .
\end{gather*}

The condition $ P_{\ell+1} \big|v_s^{(\ell)}\big\rangle  = 0 $ yields
\begin{gather*}
\sum\limits_{k=1}^{N^{(\ell)}} \sum\limits_{j=1} \alpha_k \binom{\ell+1}{j} \frac{k!}{(k-j)!} H^{k-j} P_{\ell+1-j}
P_{\ell-1}^{N^{(\ell)}-k} \big|u_0^{(\ell)}\big\rangle  = 0.
\end{gather*}
From the above relation we can see that $ \alpha_k = 0$ $(k \neq 0)$.
We also see that $ \alpha_0 = 0 $ from
\begin{gather*}
0 = C \big|v_s^{(\ell)}\big\rangle  = \alpha_0 N^{(\ell)} (\ell+1) p P_{\ell-1}^{N^{(\ell)}-1} \big|u_0^{(\ell)}\big\rangle .
\end{gather*}
Thus there are no singular vectors in $ V^{(\ell)}_{\ell}$.
That is, $ V^{(\ell)}_{\ell} $ is an irreducible ${\mathfrak g}_{\ell}$-module.
